The experimental setup, shown schematically in Fig. 1, consists of four subsystems: a 3D-printed PB frame, a droplet injection and positioning stage, a vacuum chamber with pressure control, and a high-speed imaging system.

The PB frame is a triangular prism (photosensitive resin, base edge a = 30 mm, height b = 45 mm, strut thickness ~3 mm) 3D-printed to fix a single Plateau border in space. After immersion in the surfactant solution and vertical extraction, capillary forces spontaneously form three soap films meeting at a central PB channel. The resulting PB has width w = 90 +/- 5 um, length l = 33.6 +/- 0.5 mm (measured from ImageJ micrographs), and the three films meet at the characteristic tetrahedral angle alpha = 109.5 +/- 0.5 deg [Fig. 1(b)]. This geometry remains stable for tens of minutes, far exceeding the time required for a complete experimental sequence. Compared to the random foam layer used in Ref. [19], this fixed-frame approach guarantees that the PB is centered beneath the injection needle for every trial.

The working fluid is a solution of 5.6 wt% commercial detergent, 5.6 wt% glycerol, and 88.8 wt% ultrapure water, stirred thoroughly and aged for at least 24 hours before use. The density is rho = 1.0 x 10^3 kg/m^3, the surface tension is sigma = 30.2 +/- 0.3 mN/m (pendant drop method), and the dynamic viscosity is mu = 1.2 +/- 0.1 mPa.s (rotational viscometer). The glycerol content enhances film stability without significantly altering the surface tension.

Droplets of the same solution are dispensed from a syringe pump (needle inner diameter 0.4-0.8 mm) positioned at a height h = 5-150 mm directly above the upper PB node. The equivalent spherical radius R = 1.0-2.2 mm is determined from side-view images using a custom MATLAB code that fits an ellipse to the droplet contour. The droplet impact velocity v0 (0.1-1.8 m/s, measured at the instant the droplet first contacts the soap film) is controlled by varying h and is extracted from image sequences via frame-to-frame centroid tracking. After droplet release, the liquid pool is raised electrically to immerse the frame, resetting the PB for the next trial.

To investigate the role of gas properties, the entire assembly (frame, syringe pump, liquid pool, and imaging optics) is enclosed in a vacuum chamber [Fig. 1(a)]. The ambient pressure P is varied from 0.1 to 1.0 atm using an external vacuum pump. Reducing P decreases the gas density proportionally (from 1.2 to 0.12 kg/m^3 for air) while the dynamic viscosity of air remains nearly constant (mu_a ~ 18 uPa.s), so the kinematic viscosity increases roughly tenfold across the pressure range — a critical factor for gas-film drainage dynamics [22].

The droplet-PB interaction is recorded from the side view by a high-speed camera (v711, Phantom, U.S.A.) at frame rates of 3000-10000 fps, with backlight illumination provided by a collimated LED source. A total of 199 experiments were performed, spanning the parameter space (R, v0, P) with approximately uniform coverage [see Fig. 2(a)].

For each recorded sequence, image analysis is automated through a multi-agent processing pipeline [23]. Briefly: Sisyphus monitors raw data and triggers tracking (circular Hough detection for droplet identification, frame-to-frame centroid displacement for velocity, and quadratic fitting in the time domain for acceleration). Franklin scores each case on tracking quality, radius stability, fit consistency, and center-jump diagnostics, assigning confidence bands (excellent, usable, or needs review). Eureka cross-references observed phenomena against a local literature corpus of 178 curated references and provides independent, literature-grounded observations without modifying processing parameters. Cases with a Franklin score >= 85 are promoted to a live web dashboard for real-time inspection. To serve as a control for the droplet-PB interaction, we additionally processed rigid-sphere (PP, PMMA, PS, POM) free-fall and multi-sphere experiments through the same pipeline, tracking particle trajectories rather than droplet oscillation frequencies [see Supplementary Material].